\documentclass[10pt,a4paper]{amsart}
\usepackage[margin=19mm]{geometry}
\usepackage{amsmath,amssymb,mathtools}
\usepackage[scr=rsfs,cal=euler]{mathalfa}
\usepackage{xfrac}
\usepackage[unicode,hidelinks,bookmarks=false]{hyperref}
\newcommand{\ie}{i.e.\ }
\newcommand{\cf}{cf.\ }
\newcommand{\resp}{respectively\ }
\newcommand{\Subsection}{\subsection}
\providecommand{\nobreakdash}{\nobreak\hbox{-}}
\setlength{\emergencystretch}{2em}
\multlinegap0pt
\usepackage{bm}
\usepackage{bbm}
\usepackage{dsfont}
\usepackage{xspace}
\usepackage{cancel}
\usepackage{mathtools}
\usepackage{amsthm}
\makeatletter
\@ifundefined{thm}{\newtheorem{thm}{Thm}}{}
\@ifundefined{prop}{\newtheorem{prop}[thm]{Proposition}}{}
\makeatother
\newcommand\Z{{\mathbb{Z}}}
\newcommand\ac{\acute}
\newcommand\bu{\bar{u}}
\newcommand\oo{\mathfrak{o}}
\newcommand\s{\sigma}%\s == \sigma
\title[Shalika newforms for GL(n)]{Shalika newforms for GL(n)}
\author{Takeo Okazaki}
\date{Source: arXiv:2403.04119v1 (CC BY 4.0)}
\begin{document}
\maketitle

\noindent\emph{Source and licence.} Shalika newforms for GL(n). Version arXiv:2403.04119v1 (CC BY 4.0), distributed under the Creative Commons Attribution 4.0 International licence, as recorded in the arXiv licence metadata for this exact version and shown on the abstract page https://arxiv.org/abs/2403.04119v1. Full rights notice: licenses/arxiv-2403.04119-CC-BY-4.0.txt.\par\smallskip

\noindent\textit{Editorial scope.} Counted unit: the Prop whose source label is \texttt{prop:LST} in the article (source0.tex lines 1144--1149, source label \texttt{prop:LST}), delivered together with the article's own complete original proof of it (source0.tex lines 1150--1334, one \texttt{proof} environment of 185 lines). No mathematics is written, repaired or re-derived: statement and proof are the article's own text line for line. Required context retained uncounted: none. Every definition, display and label of those lines is retained unchanged. Omitted from this package and not redistributed: 1--1143, 1335--3206. Nothing inside a retained excerpt is omitted or altered: every retained body is delivered exactly as the article's lines are written, so no omission receipt of any kind accompanies this package. Citations: the retained excerpts carry no citation command at all, so no bibliography entry of the article is transcribed and no reference is redistributed. Reference adaptation: none was needed and none was made; every reference inside the delivered unit resolves inside it, so no unresolved reference or citation is produced. Printed slips kept literal: no printed word, symbol, hypothesis or number of the counted statement or of its proof was altered. Layout and class: the article's own class is \texttt{amsart} with packages including graphicx, mathrsfs, oubraces, float, array, booktabs, enumerate, amsfonts, amsmath, amssymb, tabularx, blkarray, multirow, bigdelim; the wrapper is the lane's pinned \texttt{amsart} editorial wrapper at 10pt on A4, which restates the theorem-like environments the retained text needs and adds the article's own macro definitions that the retained text needs (\texttt{\textbackslash Z}, \texttt{\textbackslash ac}, \texttt{\textbackslash bu}, \texttt{\textbackslash oo}, \texttt{\textbackslash s}), with extra wrapper packages (bm, bbm, dsfont, xspace, cancel, mathtools). The delivered numbering is the unit's own. Assets: no retained span contains a figure environment, a graphics-inclusion command, a PDF-page inclusion, a table or a TikZ picture; no image file is copied into this package, the package declares no asset dependency and no image file is redistributed with it. Licence: version arXiv:2403.04119v1 of this article is distributed under the Creative Commons Attribution 4.0 International licence, as recorded in the arXiv licence metadata for this exact version and shown on the abstract page https://arxiv.org/abs/2403.04119v1. A full rights notice is delivered at licenses/arxiv-2403.04119-CC-BY-4.0.txt. Characters: every retained excerpt is pure ASCII, so the delivered engine, XeLaTeX with its default Latin Modern text font, reports no missing character.
\begin{prop}\label{prop:LST}
With the above notation, $\bu_x$ lies in the double coset 
\begin{align*}
\ac{H_f} u(x(\tau_x^f)^*) d(x(\tau_x^f)) K^1_{r+1}.
\end{align*}  
\end{prop}

\begin{proof}
We prove by induction on $r$.
For $r = 1$, the statement follows immediately from the the Gauss decomposition: 
\begin{align}
\begin{bmatrix}
1&  \\
y &1
\end{bmatrix} = \begin{bmatrix}
1& y^{-1}  \\
&1
\end{bmatrix}
\begin{bmatrix}
& y^{-1} \\
-y &
\end{bmatrix}\begin{bmatrix}
1&y^{-1}  \\
&1
\end{bmatrix}, \ \ y \neq 0. \label{eq:GD}
\end{align} 
Assume the induction hypothesis that the statement for $r = i$ is true.
Let $x \in F^{i+1}$, and $f \in \Z^{i+1}$.
The statement for $i+1$ is obvious if $x$ lies in $\oo^{i+1}$.
Assume $x \not\in \oo^{i+1}$.
Conjugating $\bu_x$ by $\ac{\s_x} \in G_{i+2}$, we may assume that $\s_x$ is the identity, i.e., $o(x_1) \ge \cdots \ge o(x_{i+1})$.
Write
\begin{align*}
x' = (x_{1}, \ldots, x_i), \ f' = (f_{1}, \ldots, f_i), \ y = x_{i+1}.
\end{align*}
Using (\ref{eq:GD}), we transform
\begin{align*}
\bar{u}_x &= 
\begin{bmatrix}
1_{i} & & \\
 &1 &  \\
x' & &1  \\
\end{bmatrix}
\begin{bmatrix}
1_{i} & & \\
 &1 &  \\
 & y &1  \\
\end{bmatrix} \\
&= 
\begin{bmatrix}
1_{i} & & \\
 &1 &  \\
x' & &1  \\
\end{bmatrix}
\begin{bmatrix}
1_{i} & & \\
 &1 & y^{-1}  \\
 & &1  \\
\end{bmatrix} 
\begin{bmatrix}
1_{i} & & \\
 & & y^{-1} \\
 & -y & \\
\end{bmatrix}
\begin{bmatrix}
1_{i} & & \\
 &1 & y^{-1}  \\
 & &1  \\
\end{bmatrix}.
\end{align*} 
Since $y^{-1} \in \oo$, $\bu_x$ lies in the left coset
\begin{align*}
& \begin{bmatrix}
1_{i} & & \\
 &1 &  \\
x' & &1  \\
\end{bmatrix}
\begin{bmatrix}
1_{i} & & \\
 &1 & y^{-1}  \\
 & &1  \\
\end{bmatrix} 
\begin{bmatrix}
1_{i} & & \\
 &y^{-1} &  \\
 & &y  \\
\end{bmatrix}K_{i+2}^1 \\
&= 
\begin{bmatrix}
1_{i} & & \\
 &1 & y^{-1}  \\
 & &1  \\
\end{bmatrix}
\begin{bmatrix}
1_{i} & & \\
 &y^{-1} &  \\
 & &y  \\
\end{bmatrix}
\begin{bmatrix}
1_{i} & & \\
 -x' &1 &  \\
& &1  \\
\end{bmatrix} 
\begin{bmatrix}
1_{i} & & \\
 &1 &  \\
y^{-1}x' & &1 \\
\end{bmatrix} K_{i+2}^1.
\end{align*}
Since $y^{-1}x'$ lies in $\oo^{i+1}$, this equals
\begin{align*} 
\begin{bmatrix}
1_{i} & & \\
 &1 & y^{-1}  \\
 & &1  \\
\end{bmatrix}
\begin{bmatrix}
1_{i} & & \\
 &y^{-1} &  \\
 & &y  \\
\end{bmatrix}
\begin{bmatrix}
1_{i} & & \\
 -x' &1 &  \\
& &1  \\
\end{bmatrix} K_{i+2}^1.
\end{align*} 
Put
\begin{align*}
\tau'' = \tau_{x'}^{f'} \cap \{j \in \{1, \ldots, i \} \mid o(x_j) -o(y) < f_{i+1} -f_j \}, 
\end{align*} 
so that 
\begin{align}
\tau_x^f = \tau'' \sqcup \{i+1 \} = \tau_{x''}^{f'} \sqcup \{i+1 \}. \label{eq:tauxf}
\end{align} 
The last left coset is contained in the double coset
\begin{align*}
H_f \begin{bmatrix}
1_{i} & & \\
 &1 & y^{-1}  \\
 & &1  \\
\end{bmatrix}
\begin{bmatrix}
1_{i} & & \\
 &y^{-1} &  \\
 & &y  \\
\end{bmatrix}\begin{bmatrix}
1_{i} & & \\
x'' &1 &  \\
& &1  \\
\end{bmatrix} K_{i+2}^1, 
\end{align*} 
where we write $x''$ for $x'(\tau'')$.
By the induction hypothesis, 
\begin{align*}
\bu(x'') \in H_{f'} u(x''(\tau_{x''}^{f'})^*)d(x''(\tau_{x''}^{f'})) K_{i+1}^1.
\end{align*} 
Since $\ac{H}_{f'}$ is contained in $H_f$, the last double coset is contained in 
\begin{align*}
& H_f \begin{bmatrix}
1_{i} & & \\
 &1 & y^{-1}  \\
 & &1  \\
\end{bmatrix}
\begin{bmatrix}
1_{i} & & \\
 &y^{-1} &  \\
 & &y  \\
\end{bmatrix} \ac{u}(x''(\tau_{x''}^{f'})^*)\ac{d}(x''(\tau_{x''}^{f'})) K_{i+2}^1 \\
=& 
H_f \ac{u}(y^{-1}x''(\tau_{x''}^{f'})^*) u(x(\tau_{x''}^{f'} \sqcup \{ i + 1\})^*) \begin{bmatrix}
1_{i} & & \\
 &y^{-1} &  \\
 & &y  \\
\end{bmatrix} \ac{d}(x''(\tau_{x''}^{f'})) K_{i+2}^1.
\end{align*} 
By (\ref{eq:tauxf}), this equals 
\begin{align*}
H_f \ac{u}(y^{-1}x''(\tau_{x''}^{f'})^*) u(x(\tau_x^f)^*) d(x(\tau_x^f)) K_{i+2}^1.
\end{align*}  
By (\ref{eq:tauxf}) and the definition of $\tau_x^f$,
\begin{align*}
o(x''(\tau_{x''}^{f'})_j) -o(y) = o(x(\tau_x^f)_j) -o(y) < f_{i+1} -  f_j
\end{align*} 
for all $ j \in \{1, \ldots, i \}$ such that $x_j \not\in \oo$.
By the lemma below, $u(y x''(\tau_{x''}^{f'})^*)$ lies in $H_{f'}$.
Thus the last double coset equals 
\begin{align*}
H_f u(x(\tau_x^f)^*) d(x(\tau_x^f)) K_{i+2}^1.
\end{align*}
This establishes the induction step.
\end{proof}

\end{document}
